\documentclass[11pt]{article}
\usepackage[T1]{fontenc}
\usepackage[utf8]{inputenc}
\usepackage{lmodern}
\usepackage[margin=1in]{geometry}
\usepackage{amsmath,amssymb,amsthm,mathtools}
\usepackage{microtype,enumitem,booktabs}
\usepackage[hidelinks]{hyperref}
\usepackage{xurl}
\newtheorem{theorem}{Theorem}
\newtheorem{lemma}[theorem]{Lemma}
\newtheorem{corollary}[theorem]{Corollary}
\newtheorem{proposition}[theorem]{Proposition}
\theoremstyle{remark}\newtheorem{remark}[theorem]{Remark}
\newcommand{\OPT}{\operatorname{OPT}}
\newcommand{\rhoD}{\rho_{\mathrm{det}}}
\newcommand{\rhoS}{\rho_{\mathrm{stoch}}}
\newcommand{\dd}{\,\mathrm dt}
\title{An Exact Quartet Reduction for the Two Highest Nontrivial Capacities}
\author{Phase 10 research candidate}
\date{7 October 2026}
\begin{document}
\maketitle

\begin{abstract}
For an arbitrary number $N\ge3$ of distinct ordered scalar states with
positive weights, we consider deterministic hierarchy prices at capacities
$(N-2,N-1)$. We prove that either compatible flat optima exist, or four
consecutive original states have exactly the two global flat distortion
values and a hierarchy price at least as large as that of the original
source. All weights remain their original values. Hierarchical cells are
allowed to be noncontiguous. Consequently the supremum of these prices
over all finite $N$ equals the weighted four-state $(2,3)$ supremum, for
each fixed strictly convex generator. For nonnegative nonzero polynomial
curvature of degree at most $r$, the inherited Phase~9 analytic bound gives
$1936r^2/(\log r)^2$ for $r\ge2$, uniformly in $N$. The new part is the
combinatorial restriction and lifting argument; the weighted quartet
polynomial estimate is reproduced with its provenance made explicit.
\end{abstract}

\section{Definitions and elementary facts}

Fix $N\ge3$, locations $x_1<\cdots<x_N$ in an interval $I$, and weights
$p_i>0$. The weights may sum to any positive number. Let $\varphi$ be
$C^2$ and strictly convex on $I$, and write $g=\varphi''\ge0$.
For a nonempty cell $C\subseteq\{1,\ldots,N\}$, set
\[
 p_C=\sum_{i\in C}p_i,\qquad
 m_C=p_C^{-1}\sum_{i\in C}p_ix_i,\qquad
 D(C)=\sum_{i\in C}p_i\varphi(x_i)-p_C\varphi(m_C).
\]
For a partition $P$, set $D(P)=\sum_{C\in P}D(C)$, and let
$\OPT_k$ be the minimum of $D(P)$ over partitions with at most $k$ cells.
The deterministic price at the capacities in question is
\[
 R_N=\min_{H_{N-1}\succeq H_{N-2}}
 \max\left\{
 \frac{D(H_{N-1})}{\OPT_{N-1}},
 \frac{D(H_{N-2})}{\OPT_{N-2}}\right\},
 \qquad |H_j|\le j.
\]
Refinement is denoted by $\succeq$. No contiguity restriction is imposed
on either hierarchy partition. Strict convexity, positive weights, and
distinct locations imply $\OPT_k>0$ for $k<N$, so all displayed ratios
are well defined. Multiplying every weight by one common positive factor
multiplies every distortion by that factor and leaves every price unchanged.

\begin{lemma}[Contiguous flat optima and deficit forms]
\label{lem:flat}
A contiguous flat optimum exists for every $k\le N$, and every flat
optimum fills its budget when $k<N$. In particular, a contiguous optimum
at capacity $N-1$ consists of one adjacent pair and all remaining
singletons. A contiguous optimum at capacity $N-2$ consists of either
one consecutive triple or two disjoint adjacent pairs, with all remaining
states singletons. Therefore
\[
 q:=\OPT_{N-1}=\min_{1\le i<N}D(\{i,i+1\}).
\]
\end{lemma}

\begin{proof}
The centroid identity for the Bregman divergence
$B_\varphi(x,a)=\varphi(x)-\varphi(a)-\varphi'(a)(x-a)$ is
\[
 \sum_{i\in C}p_iB_\varphi(x_i,a)
 =D(C)+p_CB_\varphi(m_C,a).
\]
Thus a weighted centroid minimizes the reporting cost of its cell. For
ordered distinct centers $a<b$, the difference
$B_\varphi(x,a)-B_\varphi(x,b)$ is affine in $x$ with positive slope
$\varphi'(b)-\varphi'(a)$. Assigning each source point to a nearest center,
with ordered tie-breaking, consequently gives contiguous cells. Starting
with centers of a flat optimum, this reassignment and subsequent centroid
updates cannot increase its distortion, and use at most $k$ cells.
Duplicate centers, if any, may first be identified. This proves existence
of a contiguous optimum.

Splitting a nonsingleton cell into its smallest-location singleton and
the other states strictly decreases its cost: their means are distinct,
so the corresponding two-group Jensen gap is strictly positive. Any
partition with fewer than $k<N$ cells can therefore be improved. Finally,
the deficit identity
$N-|P|=\sum_{C\in P}(|C|-1)$ gives the two asserted forms. Every adjacent
pair supplies a feasible $(N-1)$-cell partition, which proves the formula
for $q$.
\end{proof}

\begin{lemma}[Two adjacent costs inside a triple]
\label{lem:triple}
For three ordered states $A<B<C$, with their original positive weights,
\[
 D(ABC)\ge D(AB)+D(BC).
\]
\end{lemma}

\begin{proof}
For any cell, the Jensen kernel is
\[
 K_C(t)=\sum_{i\in C}p_i(x_i-t)_+-p_C(m_C-t)_+,
 \qquad D(C)=\int_I K_C(t)g(t)\dd.
\]
The kernel identity follows by integrating $\varphi''$ twice; affine
terms cancel. For each fixed $t$, the hinge $x\mapsto(x-t)_+$ is convex.
The Jensen gap of a union equals the sum of the within-group Jensen gaps
plus the nonnegative Jensen gap of the group centroids. Applying this to
$AB$ and $C$ gives $K_{ABC}\ge K_{AB}$; applying it to $A$ and $BC$ gives
$K_{ABC}\ge K_{BC}$. The two pair kernels have disjoint open supports,
$(A,B)$ and $(B,C)$, and vanish at the common endpoint. Hence pointwise
$K_{ABC}\ge K_{AB}+K_{BC}$. Integrate against $g\ge0$.
\end{proof}

\section{The exact restriction theorem}

\begin{theorem}[Consecutive quartet domination]
\label{thm:reduction}
Choose any cheapest adjacent pair $e$ of cost $q=\OPT_{N-1}$ and any
contiguous flat optimum $P$ at capacity $N-2$, of cost $F=\OPT_{N-2}$.
One of the following conclusions holds.
\begin{enumerate}[label=\textup{(\roman*)}]
\item There is a hierarchy with both levels flat optimal, so $R_N=1$.
\item The only nonsingleton cells of $P$ are two pairs $AB$ and $CD$ on
four consecutive original states $A<B<C<D$, and $e=BC$. For the source
restricted to $S=\{A,B,C,D\}$, keeping its original weights, its flat
values satisfy
\[
 \OPT_3(S)=q,\qquad \OPT_2(S)=F=D(AB)+D(CD),
\]
and its unrestricted deterministic $(2,3)$ hierarchy price satisfies
\[
 R_N\le R_4(S).
\]
\end{enumerate}
In particular, if $R_N>1$, conclusion \textup{(ii)} holds for every such
choice of $e$ and $P$. When $N=3$, necessarily $R_N=1$.
\end{theorem}

\begin{proof}
The partition whose sole nonsingleton cell is $e$ is flat optimal at the
fine capacity. We first check all possible forms of $P$.

Suppose the sole nonsingleton cell of $P$ is a triple $T=ABC$ of consecutive
states. If both endpoints of $e$ belong to $T$, the fine partition just
described refines $P$, proving (i). If $e$ is disjoint from $T$, take
$f=AB$ and $h=BC$. Then $e$ and $f$ are disjoint, and
\[
 D(e)+D(f)\le D(h)+D(f)\le D(T)=F,
\]
because $e$ is globally cheapest among adjacent pairs and
Lemma~\ref{lem:triple} applies. The partition with nonsingleton cells
$e,f$ has $N-2$ cells and cost at most $F$, hence exactly $F$. It is a
flat optimum compatible with the chosen fine partition. If $e$ meets
$T$ in exactly one state, adjacency of $e$ and consecutivity of $T$
force that intersection to be an endpoint of $T$. If the endpoint is
$A$, take $f=BC$ and $h=AB$; if it is $C$, take $f=AB$ and $h=BC$.
Again $e$ is disjoint from $f$, and the same inequality proves (i).
These cases exhaust the triple form.

Suppose instead that the only nonsingleton cells of $P$ are disjoint
adjacent pairs $f$ and $h$, with $F=D(f)+D(h)$. If $e$ equals either of
them, the chosen fine partition refines $P$. If $e$ is disjoint from at
least one of them, say $f$, then the partition with nonsingleton cells
$e,f$ is feasible at capacity $N-2$, and
\[
 D(e)+D(f)\le D(h)+D(f)=F.
\]
It is therefore a compatible flat optimum. Notice that this argument
allows $e$ to meet $h$, or to be disjoint from both $f$ and $h$.

The sole case remaining is that $e$ meets both $f$ and $h$. Write the two
coarse pairs in order as $AB$ and $CD$. An adjacent pair meeting both
must be $BC$. Thus $A,B,C,D$ are four consecutive original states,
and $P$ and $e$ have the form in (ii).

It remains to prove the exact flat-value identities and the price
comparison. The restricted quartet has a three-cell partition with pair
$BC$, of cost $q$, and the two-cell partition $AB\mid CD$, of cost $F$.
Conversely, any partition of $S$ with at most three (respectively two)
cells extends to a partition of the entire source with at most $N-1$
(respectively $N-2$) cells by making every state outside $S$ a singleton.
Its distortion is unchanged because every outside singleton has zero
cost. This rules out restricted values smaller than $q$ or $F$ and proves
the two identities. This step applies to arbitrary restricted partitions,
including noncontiguous ones.

Finally, apply the same singleton extension simultaneously to both levels
of any restricted hierarchy. Refinement is preserved, distortion is
unchanged, and the denominators are precisely the same. The full-source
optimum can only be better than this feasible hierarchy, proving
$R_N\le R_4(S)$. Restricted optimal hierarchies exist because their
partition set is finite; equivalently one can take the infimum of the
same comparison over all of them. For $N=3$ the coarse capacity is one,
whose single cell is compatible with every fine partition.
\end{proof}

\begin{corollary}[Equality of the all-size and quartet envelopes]
\label{cor:envelope}
For a fixed $C^2$ strictly convex generator $\varphi$ on $I$, let
$\mathcal C_N(\varphi)$ be the supremum of $R_N$ over $N$ distinct
locations in $I$ and positive weights. Then
\[
 \mathcal C_N(\varphi)\le\mathcal C_4(\varphi)\quad(N\ge3),
 \qquad
 \sup_{N\ge3}\mathcal C_N(\varphi)=\mathcal C_4(\varphi).
\]
The same equality holds after taking a supremum over any class of such
generators, in particular nonnegative nonzero polynomial curvatures of
degree at most $r$ on $[0,1]$.
\end{corollary}

\begin{proof}
Every hierarchy price is at least one, and the quartet envelope is at
least one. Apply Theorem~\ref{thm:reduction} pointwise. Taking the
supremum over $N$ includes $N=4$, so the reverse inequality for the last
display is automatic. If probabilities summing to one are required,
divide every weight in the restricted quartet by the same total mass;
the price is unchanged by this common scaling. No individual cell is
conditionally renormalized.
\end{proof}

\begin{remark}[Meaning and limits of ``exact'']
The reduction has no multiplicative loss and preserves both flat
denominators exactly before any common weight normalization. It asserts
domination by a restriction, not equality of an individual source price
with the selected quartet price. The lifted quartet hierarchies form
only a subset of the global feasible hierarchies. The corollary asserts
equality of the supremum over all finite source counts with the quartet
supremum; it does not prove equality separately for every fixed $N>4$.
\end{remark}

\section{The polynomial consequence}

\begin{theorem}[Uniform degree bound at the top two capacities]
\label{thm:degree}
Let $N\ge3$, let $x_1<\cdots<x_N$ lie in $[0,1]$ with arbitrary positive
weights, and suppose $g=\varphi''$ is a nonzero polynomial of degree at
most $r$, nonnegative on $[0,1]$. For the deterministic price $R_N$ at
capacities $(N-2,N-1)$:
\begin{align*}
 R_N\ge128,\ r\ge1
 &\quad\Longrightarrow\quad
 R_N^2\le16r^2\exp\left(\frac{20r}{\sqrt{R_N}}\right),\\
 R_N\ge256,\ r\ge1
 &\quad\Longrightarrow\quad
 R_N(\log R_N)^2\le1936r^2.
\end{align*}
For every integer $r\ge2$,
\[
 \boxed{\displaystyle R_N\le1936\frac{r^2}{(\log r)^2}.}
\]
The constant is independent of the source count, all positive weights,
and all positive posterior gaps. The inherited Phase~9 matching quartet
lower construction makes the supremum over all $N\ge3$ have sharp order
$\Theta(r^2/(\log r)^2)$.
\end{theorem}

\begin{proof}
A nonzero nonnegative polynomial vanishes on no interval, so its
twice-antiderivative is strictly convex. Apply
Theorem~\ref{thm:reduction}. There is nothing to prove if $R_N=1$;
otherwise $R_N\le R_4(S)$ for the specified quartet. The weighted quartet
bound, proved in Phase~9 and reproduced in
Proposition~\ref{prop:inherited}, gives the final displayed estimate.
For the implicit inequality, the quartet estimate can be written
\[
 R_4(S)^2\exp\left(-\frac{20r}{\sqrt{R_4(S)}}\right)\le16r^2.
\]
The function $s\mapsto s^2\exp(-20r/\sqrt{s})$ is strictly increasing
for $s>0$, so this inequality transfers to $R_N$ whenever $R_N\ge128$.
Likewise $s\mapsto s(\log s)^2$ is increasing for $s\ge1$, which gives
the second implication. The assertion about the envelope's lower order
uses the stated, inherited quartet witnesses and the fact that the
all-size supremum includes $N=4$; no padding argument is used.
\end{proof}

\begin{remark}[Stochastic upper bound only]
Every deterministic hierarchy is an admissible stochastic hierarchy,
with the same flat denominators. Thus the same uniform upper bound
holds for the stochastic hierarchy price. The restriction proof itself
is a deterministic argument: it does not establish an equality between
the all-size stochastic envelope and the stochastic quartet envelope.
The inherited stochastic quartet lower witnesses still establish the
matching all-size stochastic degree order.
\end{remark}

\section{Inherited weighted quartet analytic input}

The following proposition, including its constants and proof strategy,
is the finalized Phase~9 weighted quartet theorem. It is reproduced here
so the upper-bound argument can be read without another file. It is not
a new Phase~10 analytic estimate. The source is
\path{phase9_research_work/research/degree_gap/weighted_four_state_upper.tex},
also included in the Phase~9
working paper \emph{Sharp Weighted Four-State Hierarchy Prices for
Polynomial Curvature}.

\begin{proposition}[Phase 9 weighted quartet estimate]
\label{prop:inherited}
For four distinct weighted states and nonnegative nonzero polynomial
curvature of degree at most $r\ge1$, let $R$ be the unrestricted
deterministic hierarchy price at capacities $(2,3)$. If $R\ge128$, then
\[
 R^2\le16r^2e^{20r/\sqrt R}.
\]
If $R\ge256$, then $R(\log R)^2\le1936r^2$. In particular, for $r\ge2$,
$R\le1936r^2/(\log r)^2$.
\end{proposition}

\begin{proof}
We record two classical polynomial estimates. If $f\ge0$ has degree at
most $r$ on an interval $E$ of length $L$, and $M_E=\max_Ef$, the Markov
inequality $\|f'\|_\infty\le2r^2M_E/L$ implies
\begin{equation}
 \int_E f\ge\frac{LM_E}{8r^2}.
 \label{eq:markov}
\end{equation}
Indeed a one-sided segment of length $L/(4r^2)$ is available from a
maximizer, and $f\ge M_E/2$ there. If $t_0$ lies at distance $d>0$ to
the right of $E$, Chebyshev extremality gives
\begin{equation}
 |f(t_0)|\le M_E T_r(1+2d/L)
 \le M_Ee^{2r\sqrt{d/L}}.
 \label{eq:cheb}
\end{equation}
For completeness, after rescaling $E$ to $[-1,1]$, any violation of the
first inequality would yield $\lambda\in(0,1)$ such that
$T_r-\lambda f/M_E$ has $r$ roots between the alternating extrema of
$T_r$ and an additional root at the external point. This contradicts
its degree. The exponential bound follows from
$T_r(\cosh s)=\cosh(rs)$ and
$\operatorname{arcosh}(1+x)\le\sqrt{2x}$.

Since $R>1$, the argument of Theorem~\ref{thm:reduction}, specialized to
four states, permits the notation
\[
 A<B<C<D,\quad q=D(BC)=\OPT_3,\quad
 Q=D(AB)+D(CD)=\OPT_2.
\]
Write $p_L=D(AB)$, $p_R=D(CD)$, $T_L=D(ABC)$, and $T_R=D(BCD)$.
Three feasible deterministic hierarchies give
\begin{equation}
 R\le\frac{\min(p_L,p_R)}q,\qquad
 R\le\frac{T_L}Q,\qquad R\le\frac{T_R}Q.
 \label{eq:candidates}
\end{equation}
For the first retain $AB\mid CD$ and merge the cheaper outer pair at
the fine level. For the others retain fine pair $BC$ and merge the
corresponding adjacent triple at the coarse level. All right sides
are at least one by the defining flat optimality. No enumeration of
noncontiguous optimal hierarchies is needed.

Reflect the source if needed so that the source weights satisfy
$p_B\le p_C$. Normalize only the coordinate of the triple $BCD$ to
\[
 B=0,\qquad C=c\in(0,1),\qquad D=1,
 \qquad (p_B,p_C,p_D)=(u,v,w),\quad 0<u\le v.
\]
The transformed curvature is $L_0^2g(B+L_0t)$, where $L_0$ is the
original triple hull length; its degree and nonnegativity are preserved.
Weights are unchanged. Denote that curvature again by $g$. The triple,
adjacent-pair, and outer-pair kernels are
\begin{align*}
 K(t)&=\min\{ut+v(t-c)_+,\ w(1-t)+v(c-t)_+\},\\
 k_q(t)&=\min\{ut,v(c-t)\}\mathbf1_{[0,c]}(t),\\
 k_p(t)&=\min\{v(t-c),w(1-t)\}\mathbf1_{[c,1]}(t),\\
 H(t)&=\min\{ut,w(1-t)\}.
\end{align*}
Put $T=T_R=\int Kg$ and $S=p_R+q=\int(k_q+k_p)g$. From
\eqref{eq:candidates},
\begin{equation}
 q\le\frac{Q}{2R}\le\frac{T}{2R^2},\qquad
 S\le Q+q\le\frac{2T}{R}.
 \label{eq:small}
\end{equation}
Two pointwise kernel comparisons control degenerating weights:
\begin{equation}
 0\le K-H\le k_q+k_p,
 \label{eq:H}
\end{equation}
and, in their respective open intervals,
\begin{equation}
 \frac{k_q(t)}{K(t)}\ge\min\{1,(c-t)/t\}\quad(t<c),
 \qquad
 \frac{k_p(t)}{K(t)}\ge\frac{t-c}{c+2(t-c)}\quad(t>c).
 \label{eq:ratios}
\end{equation}
For $t<c$, subtracting $H$ from $K$ amounts to adding $v(c-t)$ to one
entry of a minimum, so the increment is between zero and
$\min\{ut,v(c-t)\}=k_q$. For $t>c$, adding $v(t-c)$ to the other entry
gives the bound by $k_p$. For the first ratio, use $K\le ut$ and
$v\ge u$. For the second, write $\Delta=t-c>0$ and use
\[
 \frac{\min\{B,C\}}{\min\{A+B,C\}}\ge\frac{B}{A+B},
 \quad A=ut,\ B=v\Delta,\ C=w(1-t),
\]
followed by $v\ge u$. Endpoint claims follow by continuity.

Assume $R\ge128$ and set
$\varepsilon=16/R\le1/8$ and
$J=\{t\in[0,1]:|t-c|\le\varepsilon c\}$. Outside $J$,
\eqref{eq:ratios} yields
$k_q+k_p\ge\varepsilon K/(1+2\varepsilon)$. Consequently,
\[
 \int_{[0,1]\setminus J}Kg
 \le\frac{1+2\varepsilon}{\varepsilon}S
 \le\frac{5T}{32},\qquad
 \int_JHg\ge T-\frac{5T}{32}-S\ge\frac{3T}{4}.
\]
Here $S\le2T/R\le T/64$. Since
$H(t)\le u(1+\varepsilon)c$ on $J$ and $|J|\le2\varepsilon c$,
there is $t_0\in J$ such that
\[
 g(t_0)\ge\frac{T}{4\varepsilon u c^2}.
\]
Consider $E=[c/4,c(1-2\varepsilon)]$. It has length at least $c/2$;
$t_0$ is to its right at distance at most $3\varepsilon c$; and
$k_q\ge2u\varepsilon c$ on $E$. Using \eqref{eq:cheb},
\[
 M_E\ge\frac{T}{4\varepsilon u c^2}
 e^{-\sqrt{24}\,r\sqrt\varepsilon}
 \ge\frac{T}{4\varepsilon u c^2}e^{-20r/\sqrt R}.
\]
Equation~\eqref{eq:markov} now implies
\[
 q\ge2u\varepsilon c\int_Eg
 \ge\frac{u\varepsilon c^2M_E}{8r^2}
 \ge\frac{T}{32r^2}e^{-20r/\sqrt R}.
\]
Together with \eqref{eq:small}, this proves the implicit inequality.

Set $z=r/\sqrt R$. The inequality implies
\[
 \log(R/16)\le20z+2\log z\le22z.
\]
For $R\ge256$, the left side is at least $\tfrac12\log R$, yielding
$\sqrt R\log R\le44r$, hence $R(\log R)^2\le1936r^2$.
For $r\ge2$ and $R\ge256$ with $R>r$, use $\log R\ge\log r$.
If $R\le r$, use $(\log r)^2\le r$; if $R<256$, use
$(\log r)^2\le r^2$ and $256<1936$. These cases give the explicit bound.
\end{proof}

\section{Provenance and scope}

The centroid identity, nearest-center contiguity for scalar Bregman
costs, Jensen kernels, Markov's inequality, and Chebyshev extremality
are classical ingredients. The positive-weight adjacent triple kernel
comparison and the complete weighted quartet analytic argument above
were already present in Phase~9. The proposed additional result here is
Theorem~\ref{thm:reduction}: the deficit-two classification and cheapest
pair exchange reduce arbitrary source counts at these capacities to a
four-state restriction with both denominators preserved. The exact
envelope identity and uniform degree corollary then follow.

These statements concern only capacities $(N-2,N-1)$. They do not improve
the degree bound for arbitrary prescribed capacities, for one hierarchy
covering all capacities, or for neural training algorithms. No optimal
leading constant or historical priority is claimed. Distinct states and
positive weights are used to ensure nonzero denominators; repeated states,
zero weights, and identically zero curvature would require a separate
convention for zero-denominator ratios and are outside the stated theorem.

\end{document}
